\documentclass[10pt]{article}
\usepackage[margin=0.9in]{geometry}
\usepackage{amsmath,amsthm,booktabs}
\usepackage[hidelinks]{hyperref}
\DeclareMathSizes{10}{10}{7}{7}
\DeclareMathSizes{9}{10}{7}{7}
\newtheorem{theorem}{Theorem}
\newtheorem{lemma}[theorem]{Lemma}
\title{Deletion and forward-distance bounds near the sharp\\Caccetta--H\"aggkvist circulant}
\author{Claude Make No Mistakes}
\date{26 November 2026}
\begin{document}
\maketitle

\begin{abstract}
Let the base oriented graph on \(3r\) cyclically labeled vertices contain
the forward arcs of lengths \(1,\ldots,r-1\). We show that any oriented
graph of minimum outdegree at least \(r\) with no directed triangle must delete
at least \(r\) of these arcs. We also give a finite SAT certificate that
for every \(n\ge15\), the minimum outdegree \(\lceil n/3\rceil\) forces a
directed triangle if the maximum forward arc length is
\(\lceil n/3\rceil+2\) or less in some cyclic labeling. The certificate uses
15 vertices and is independent of \(n\).
\end{abstract}

For an arc \(u\to v\) in a given cyclic labeling of \(n\) vertices, its
forward length is the representative of \(v-u\bmod n\) in
\(\{1,\ldots,n-1\}\). An oriented graph has no loops or reciprocal arc pairs.
The triangle case of the Caccetta--H\"aggkvist conjecture asks whether every
oriented graph of order \(n\) and minimum outdegree at least \(n/3\) has a
directed triangle. Hladk\'y, Kr\'al', and Norin prove the threshold
\(0.3465n\) \cite{hkn}. At \(n=3r\), the graph with all forward lengths
\(1,\ldots,r-1\) has minimum outdegree \(r-1\) and no directed triangle,
since three such lengths sum to less than \(3r\). We call it the base graph.
The results below use the geometry of that cyclic labeling.

\section{The deletion bound}

\begin{theorem}
Let \(r\ge3\) and let \(D\) be an oriented graph on the base vertex set
with minimum outdegree at least \(r\). If \(D\) has no directed triangle,
then at least \(r\) base arcs are absent from \(D\).
\end{theorem}

\begin{proof}
Write \(F\) for the absent base arcs and assume \(f=|F|\le r-1\).
Call a vertex clean if no arc of \(F\) enters it. All vertex arithmetic
below is modulo \(3r\). Every vertex has at least one new outgoing arc,
where a new arc has length at least \(r\). A clean vertex has new arcs
of length at most \(2r\), since any longer arc would reciprocate a
retained incoming base arc.

Suppose a new arc \(u\to v\) from a clean vertex \(u\) has length \(r+a\),
with \(1\le a\le r\). For \(1\le j<a\), there is a base return path
\[
   v\ \longrightarrow\ u+2r+j\ \longrightarrow\ u .
\]
The last arc is retained, so every first arc is in \(F\). These \(a-1\)
distinct deletions leave \(v\) in need of at least \(a\) new outgoing arcs.
In particular, if both \(u\) and \(v\) are clean and \(v\to x\) is new of
length \(r+b\), the sum \(2r+a+b\) must exceed \(3r\). A smaller sum closes
a triangle with a retained base arc into \(u\), and equality makes the
new arcs reciprocal. The two new arcs force at least \(a-1\) deletions
leaving \(v\) and at least \(b-1\) leaving \(x\). Hence
\(r-1\ge f\ge a+b-2\ge r-1\), so \(b=r+1-a\) is the only available new
length parameter at \(v\). This is impossible for \(a\ge2\), since \(v\)
then needs at least two new outgoing arcs. If \(a=1\), then \(b=r\) and
all deletions leave \(x\), which makes \(x\) clean. Applying the already
excluded \(a\ge2\) case to \(v\to x\) gives a contradiction. Thus every
new arc between clean vertices has length \(r\).

In fact no new arc anywhere has length \(2r\). If \(w\to x\) has this
length, its \(r-1\) edge-disjoint base return paths
\(x\to x+j\to w\) for \(1\le j\le r-1\) must each lose an arc. This uses
all \(f\) deletions and confines their heads to
\(\{x+1,\ldots,x+r\}\). In particular \(x\) is clean. A new step \(r\)
from \(x\) reciprocates \(w\to x\), and every longer possible new step
from \(x\) has a clean head. Thus \(x\) would have no new outgoing arc.

Call a residue class modulo \(r\) clean if all three of its vertices
are clean. There is at least one such class. Partition the residue
cycle into gaps, each starting with a clean class and followed by
the \(k\ge0\) nonclean classes up to the next clean class. If each gap
with \(k\) nonclean classes had at least \(k+1\) incoming deletions,
then \(f\ge r\). Hence some gap has exactly one deletion entering
each of its \(k\) nonclean classes. Number these classes
\(1,\ldots,k\) relative to its initial clean class.

The three vertices in the initial clean class would form a triangle
if all took step \(r\). Choose one, denoted \(u\), missing this step.
It has a new arc \(u\to v\) of length \(r+t\), with \(1\le t\le r-1\).
If \(t\ge k+2\), the forced deletion from \(v\) into \(u+2r+k+1\)
enters the next clean class. If \(t=k+1\), the new arc has a clean
head. Both are impossible. Thus \(1\le t\le k\), and \(v\) is the
unique head of a deleted arc in class \(t\). The forced deletions
above show that \(v\) needs at least \(t\) new outgoing arcs.

For a new arc \(v\to x\) of length \(e\le2r-1\), the sum \(r+t+e\)
must exceed \(3r\), by the surviving base arcs into \(u\) and the
absence of reciprocal arcs. So the possible lengths lie in
\(\{2r-t+1,\ldots,2r-1\}\), a set of size \(t-1\).
Each longer new arc must reverse a deleted incoming base arc at \(v\),
since step \(2r\) is impossible. There is exactly one possible
reversal. To get \(t\) new arcs, \(v\) must take it and all \(t-1\)
possible shorter lengths.

If \(t=1\), the sole deletion into \(v\) is \(y\to v\) and the sole new
arc at \(v\) reverses it. If the deleted length were at least two,
the base return path \(y\to v-1\to v\) would survive. Its first arc
enters the initial clean class and its second is distinct from the
only deletion entering class 1. Thus \(y=v-1\), and \(y\) now needs
two new outgoing arcs. Step \(r+1\) from \(y\) would have a clean
head \(v+r\), and any larger new step would force a deletion into
\(y+2r+1=v+2r\). That vertex is also clean. Consequently \(y\) can
take only step \(r\), a contradiction.

If \(t\ge2\), the forced step \(2r-1\) from \(v\) lands at \(u+t-1\).
Its base return path through \(u+r\) has lengths \(r-t+1\) and \(t\).
The first arc enters the initial clean class, so \((u+r)\to v\)
must be the sole deleted arc entering \(v\). Its reversal forms the
triangle \(v\to u+r\to u+r+1\to v\). The middle arc survives because
the only deletion into class 1 is \(v\to u+2r+1\), already forced by
\(t\ge2\). The last arc survives since it differs from \((u+r)\to v\).
This exhausts the cases and proves \(f\ge r\).
\end{proof}

At order \(153\), the bound gives at least 51 deletions in every cyclic
labeling. The degree condition also forces at least \(153+f\) added
arcs, for at least 255 changes to directed arc membership.

\begin{lemma}
For \(r\ge3\), an oriented graph of minimum outdegree at least \(r\)
that retains all base arcs contains at least \(r\) directed triangles.
This bound is attained.
\end{lemma}

\begin{proof}
Keep the base arcs and one selected new arc per vertex. For a residue
class with all three selected lengths equal to \(r\), use its triangle.
In every other class choose a vertex \(u\) with selected length \(d>r\).
If \(d\ge r+2\), the arc has a two-edge base return path. If \(d=r+1\),
follow the selected arc at its head \(v\), of length \(e\). For
\(e\le2r-2\) there is a closing base arc into \(u\). Length \(2r-1\)
would reciprocate the selected arc from \(u\). The only remaining
option is \(e=2r\), as longer arcs reciprocate base arcs. In this
case use one of the \(r-1\) two-edge base return paths for the selected
arc at \(v\).

Triangles with one selected arc identify that arc uniquely. They are
needed at most twice per selected arc, once for its own class and once
for the unique vertex at forward distance \(r+1\) behind its tail.
In the latter case the selected length is \(2r\), so \(r-1\ge2\)
paths suffice. Each triangle with two selected arcs identifies its
initial \(u\) as the head of the sole base arc. Triangles with three
selected arcs come from distinct classes. Thus the classes give at
least \(r\) distinct triangles. Adding all arcs of length \(r\) to
the base graph gives exactly \(r\).
\end{proof}

\section{The forward-distance bound}

\begin{theorem}
Let \(n\ge15\) and \(r=\lceil n/3\rceil\). If an oriented graph on
\(n\) vertices has minimum outdegree at least \(r\) and admits a cyclic
labeling in which every arc has forward length at most \(r+2\), then
it contains a directed triangle.
\end{theorem}

\begin{proof}
Choose three cyclic gaps \(g_0,g_1,g_2\) with \(5\le g_j\le r\) and
\(\sum g_j=n\), which exist since \(15\le n\le3r\). Take five consecutive
vertices at the start of each gap. Label them \((i,j)\) with
\(0\le i\le4\), where \(j\) is the gap index modulo three. For each
selected vertex in gap \(j\), consider the seven possible outgoing
arcs of lengths \(g_j+d\) for \(-4\le d\le2\). These are distinct legal
lengths in \([1,r+2]\). At most two of these arcs can be absent since
there are only \(r+2\) possible outgoing lengths in the entire graph.

Let \(x_{i,j,d}\) indicate that this arc is absent. If \(i+d\in[0,4]\),
the arc ends at \((i+d,j+1)\). For every \(i_0,i_1,i_2\in[0,4]\) whose
cyclic differences \(d_j=i_{j+1}-i_j\) lie in \([-4,2]\), a triangle-free
graph satisfies
\[
    x_{i_0,0,d_0}\ \lor\ x_{i_1,1,d_1}\ \lor\ x_{i_2,2,d_2}.
\]
The omission limit gives a negative clause for each triple of variables
in a row. There are 105 variables, \(15\binom73=525\) row clauses and
80 triangle clauses. The archived DRAT refutation of these 605 clauses
is checked by DRAT-trim \cite{drat}. Thus the hypothesized graph
cannot be triangle-free.
\end{proof}

The same cyclic labeling gives a short direct proof for \(n=3r\) and
maximum forward length \(r+1\), with \(r\ge3\). For a triangle-free
graph, every translate of the directed triangle with ordered lengths
\((r-1,r,r+1)\) must lose an arc. Counting across translates, the total
number of omissions among those three lengths is at least \(n\).
Each vertex has at most one omission in all \(r+1\) lengths, so equality
holds, with no omissions at length \(r-2\). Every ordered permutation
of the mixed lengths therefore loses exactly one arc per translate.
An omission of step \(r\) at \(u\) forces both steps \(r-1,r+1\)
to survive at \(u+r\), by using the two mixed triangles with step \(r\)
first. Hence step \(r\) is also omitted at \(u+r\). In each residue
triple at least one step \(r\) must be omitted. Propagation forces all
steps \(r\) to be omitted, so every cycle with lengths
\((r-2,r+1,r+1)\) survives.

\section{Search at order 153}

The public SAT project reports closing the orders through 152 and lists
153 as its first remaining order \cite{campaign}. To search directly,
we encoded all allowed arcs for \(n=153\) with maximum forward length
\(51+s\), starting with \(s=0,1,2\). Here any directed triangle in
the host has total length 153 and uses only steps at least \(51-2s\).
Thus all shorter arcs can be kept. The encoding allows at most \(s\)
omitted arcs per row among the remaining lengths. The order-153
formula for \(s=2\) has 2,601 variables and 4,947 clauses and a separate
checked refutation. The wider cases \(s=3,4\) remained unresolved at
300,000 and 100,000 conflicts respectively.

We also tested 2.4 million swap proposals across eight seeds, maintaining
outdegree 51 and orientation at each accepted state. The forward-length
caps for ordinary moves ranged from 52 to 128, with longer steps allowed
in six runs. A move can replace an arc in a single row, or replace two
arcs to repair a reciprocal pair. The starting graph with steps
\(1,\ldots,51\) has exactly 51 directed triangles, one in each residue
class modulo 51. The best genuinely changed graph visited had 52.
All recorded improvements were checked by direct triangle enumeration.
This search produced no counterexample.

For context, Hamidoune's theorem covers vertex-symmetric graphs
\cite{ham}, and circular interval digraph results use constraints on
the intervals spanned by arcs \cite{css}. The deletion theorem here
allows arbitrary new arcs and the forward-distance theorem allows any
subset of the arcs in the prescribed band. Searches of indexed
literature and public X results located partial results, including a
fixed-outdegree-six announcement \cite{six}. The general triangle
case remains open.

\begin{thebibliography}{9}
\bibitem{hkn}
J.~Hladk\'y, D.~Kr\'al', and S.~Norin, \emph{Counting flags in
triangle-free digraphs}, \href{https://arxiv.org/abs/0908.2791}{arXiv:0908.2791}.
\bibitem{drat}
N.~Wetzler, M.~J.~H.~Heule, and W.~A.~Hunt Jr.,
\emph{DRAT-trim},
\href{https://doi.org/10.1007/978-3-319-09284-3_31}{SAT 2014, 422--429}.
\bibitem{campaign}
\emph{Attack log for Caccetta--H\"aggkvist directed triangles},
\href{https://github.com/wustep/maths/blob/main/problems/caccetta-haggkvist-k3/ATTACK.md}{wustep/maths}.
\bibitem{ham}
Y.~O.~Hamidoune, \emph{On iterated image size for point-symmetric relations},
\href{https://arxiv.org/abs/0704.0459}{arXiv:0704.0459}.
\bibitem{css}
M.~Chudnovsky, P.~Seymour, and B.~Sullivan, \emph{Cycles in dense digraphs},
\href{https://web.math.princeton.edu/~mchudnov/kover2.pdf}{author manuscript}.
\bibitem{six}
Fixed-outdegree-six announcement, \href{https://x.com/jjpcodes/status/2079543801387815133}{X, 21 July 2026}.
\end{thebibliography}
\end{document}
